\rhead{CX202: Design \& Analysis of Algo. Lab}
\phantomsection 
\section*{Design \& Analysis of Algorithms (CS202)}
\label{course:CX202}

\noindent \small{\hyperref[main:scheme]{$\leftarrow$ Back to Scheme}} \vspace{0.5cm}

\noindent \textbf{Credits:} 2 (0-0-2)

\subsection*{Course Content}
\noindent\textbf{Lab 1: Empirical Analysis of Sorting Algorithms} \
Focuses on experimentally verifying theoretical time complexities. Students will implement and compare the runtime performance of an O($n^2$) sort (e.g., Insertion Sort) with an O($n \cdot log n$) sort (e.g., Merge Sort) across various input sizes. \\

\vspace{0.2cm}

\noindent\textbf{Lab 2: Divide and Conquer (Quicksort)} \
Focuses on the practical implementation of the Divide and Conquer paradigm. Students will implement the Quicksort algorithm, paying close attention to the partitioning logic and the choice of the pivot element. \\

\vspace{0.2cm}

\noindent\textbf{Lab 3: Greedy Algorithms (Knapsack and Job Sequencing)} \
Focuses on making locally optimal choices to achieve a global optimum. Students will implement the greedy solution for the Fractional Knapsack problem and the Job Sequencing with Deadlines problem. \\

\vspace{0.2cm}

\noindent\textbf{Lab 4: Minimum Spanning Trees} \
Focuses on applying greedy algorithms to graph structures. Students will implement either Prim's algorithm (using a priority queue) or Kruskal's algorithm (using a Disjoint Set Union data structure) to find the MST of a given graph. \\

\vspace{0.2cm}

\noindent\textbf{Lab 5: Dynamic Programming (Memoization)} \
Focuses on the top-down approach to solving problems with overlapping subproblems. Students will implement a recursive solution to the 0/1 Knapsack problem, using a memoization table to store and reuse results. \\

\vspace{0.2cm}

\noindent\textbf{Lab 6: Dynamic Programming (Tabulation)} \
Focuses on the bottom-up approach to dynamic programming. Students will implement an iterative solution for the Longest Common Subsequence (LCS) problem using a 2D table to build up the solution from smaller subproblems. \\

\vspace{0.2cm}

\noindent\textbf{Lab 7: Backtracking (N-Queens)} \
Focuses on solving constraint satisfaction problems by exploring a state-space tree. Students will implement a backtracking algorithm to find all possible solutions for placing N queens on an N×N chessboard without attacking each other. \\

\vspace{0.2cm}

\noindent\textbf{Lab 8: Shortest Path Algorithms (Dijkstra's)} \
Focuses on finding the most efficient path in a network. Students will implement Dijkstra's algorithm to find the single-source shortest paths in a weighted graph with non-negative edge weights, using a priority queue for efficiency. \\

\vspace{0.2cm}

\noindent\textbf{Lab 9: String Matching (KMP Algorithm)} \
Focuses on efficient pattern searching within text. Students will implement the Knuth-Morris-Pratt (KMP) algorithm, which involves pre-processing the pattern to create a Longest Prefix Suffix (LPS) array to avoid redundant comparisons. \\

\vspace{0.2cm}

\noindent\textbf{Lab 10: Shortest Path with Negative Weights (Bellman-Ford)} \
Focuses on handling more complex graph scenarios. Students will implement the Bellman-Ford algorithm, which can compute single-source shortest paths in weighted graphs that may contain negative edge weights, and also detect negative cycles.

\newpage
\rhead{Curriculum Schema}